\documentclass{article}
\usepackage{amsmath}
\usepackage{amssymb}
\usepackage{amsthm}
\usepackage{mathtools}
\usepackage{geometry}
\usepackage{hyperref}

\newcommand{\rk}{\operatorname{rank}}
\newcommand{\GL}{\operatorname{GL}}
\newcommand{\R}{\mathbb{R}}
\newcommand{\Mat}{\mathrm{Mat}}
\newcommand{\Stab}{\operatorname{Stab}}

\theoremstyle{definition}
\newtheorem{theorem}{Theorem}[section]
\newtheorem{proposition}[theorem]{Proposition}
\newtheorem{lemma}[theorem]{Lemma}
\newtheorem{corollary}[theorem]{Corollary}
\newtheorem{definition}[theorem]{Definition}
\newtheorem{remark}[theorem]{Remark}

\title{Stratification Consequence}
\date{}

\begin{document}
\maketitle

\section{Stratification Consequence of the First Conjecture}

Assume the orbit-bijection conjecture from \texttt{first\_conjecture\_draft.tex}.

\begin{proposition}[Formal stratification consequence]
Let
\[
\Gamma_{d,S}^r
=
\bigsqcup_{\alpha\in A}\mathcal O_\alpha
\]
be a decomposition into locally closed $G_d^{W_0}$-stable strata. Assume that each stratum is a union of $G_d^{W_0}$-orbits and that the indexing set $A$ is finite. Let
\[
H_S:=\Stab_{G_d^{W_0}}(P_SW^\ast).
\]

Then the critical locus admits the induced decomposition
\[
\mathcal C_{d,S}^r
=
\bigsqcup_{\alpha\in A}\mathcal C_\alpha,
\qquad
\mathcal C_\alpha := \mathcal C_{d,S}^r \cap \mathcal O_\alpha,
\]
and this decomposition is a stratification by locally closed $H_S$-stable pieces.
\end{proposition}

\begin{proof}[Formal argument]
By the orbit-bijection conjecture, every $G_d^{W_0}$-orbit in $\Gamma_{d,S}^r$ meets $\mathcal C_{d,S}^r$, and its intersection with the critical locus is exactly one $H_S$-orbit.

Therefore the $G_d^{W_0}$-orbit partition of $\Gamma_{d,S}^r$ induces an
$H_S$-orbit partition of $\mathcal C_{d,S}^r$. Grouping the critical
$H_S$-orbits according to the ambient strata $\mathcal O_\alpha$ gives the
decomposition above. Since each $\mathcal O_\alpha$ is locally closed and
$H_S\subseteq G_d^{W_0}$, each $\mathcal C_\alpha$ is locally closed and
$H_S$-stable.
\end{proof}

\begin{corollary}[If the cyclic strata are indexed by combinatorial data]
Suppose the strata of $\Gamma_{d,S}^r$ are indexed by cyclic rank patterns, or by any equivalent finite combinatorial invariant:
\[
\Gamma_{d,S}^r
=
\bigsqcup_{\rho\in \mathcal R_{d,S}^{r,\mathrm{cyc}}}\mathcal O_\rho.
\]
Then the critical locus inherits the corresponding combinatorial stratification
\[
\mathcal C_{d,S}^r
=
\bigsqcup_{\rho\in \mathcal R_{d,S}^{r,\mathrm{cyc}}}\mathcal C_\rho,
\qquad
\mathcal C_\rho := \mathcal C_{d,S}^r \cap \mathcal O_\rho.
\]
\end{corollary}

\begin{remark}
This is the precise role of the first conjecture in the overall paper: it reduces the stratification problem for the critical locus to the stratification problem for the weight-space cyclic locus.
\end{remark}

\begin{remark}[What remains to be proved later]
The proposition above is only formal once the orbit-bijection conjecture is granted. The real mathematical work is therefore pushed into two later steps:
\begin{itemize}
\item classify the $G_d^{W_0}$-orbits in $\Gamma_{d,S}^r$;
\item prove that every such orbit contains exactly the right critical representative.
\end{itemize}
\end{remark}

\end{document}
